% Drafting owner: Andrea. Keep scientific content student-authored.
\subsection{Research Design and Methodological Rationale}

We propose a controlled, quantitative 2D simulation of two unmanned surface vehicles (USVs) surveying a known area under intermittent direct radio communication. We compare reactive online resynchronization with prediction-aware pre-outage task reservation; a fixed offline partition is a communication-independent reference. Channel-fade persistence (target mean fade duration) is the principal varying condition. Mission completion time is the primary outcome, with redundant coverage and communication cost as secondary measures.

The two online strategies will share the motion and sensing models, task allocator, pre-generated fade and slowdown traces, update interval, and maximum message budget; only the outage-handling policy will differ. Paired simulations can therefore test whether anticipating an outage changes allocations and outcomes. Real-boat experiments could strengthen external validity, but are infeasible within the available resources and make link conditions harder to control. A single distance-loss sensitivity check will probe the simplified model. Software tests verify implementation, not maritime realism; conclusions will be limited to the modeled two-USV scenario rather than real-world navigation safety, radio performance, or seabed-map accuracy.
